\documentclass[11pt,letterpaper]{article}
\usepackage{cornell-notes}
\usepackage{needspace}

\newcommand{\CornellDocumentTitle}{Chapter 93: an overview of cyber attacks and defenses on intelligent connected vehicles}
\title{\CornellDocumentTitle}
\author{}
\date{}
\setDocTitle{\CornellDocumentTitle}
\setDocSubtitle{Cybersecurity Cornell Notes}
\setCornellCollection{Cybersecurity Reference Notes}
\setCornellUnitType{Chapter}
\setCornellUnitNumber{93}
\setCornellUnitTitle{an overview of cyber attacks and defenses on intelligent connected vehicles}
\setCornellCompanionLabel{Cybersecurity study companion}

\begin{document}
\maketitle
\begin{CornellOverviewBox}{Chapter Orientation}
\textbf{Chapter orientation.} This chapter examines an overview of cyber attacks and defenses on intelligent connected vehicles within the broader domain of connected and automated vehicle security. Its central problem is how to reason about avs, attacks, communication, and ota while keeping security objectives, operating assumptions, evidence, and recovery connected. A practitioner should approach the material through physical consequences, safety, availability, environmental dependencies, remote connectivity, maintenance, and recovery. Useful prerequisites include basic risk, threat, control, and systems concepts.
\end{CornellOverviewBox}
\section*{Learning Objectives}
\begin{itemize}
\item Explain the security purpose and scope of An Overview of Cyber Attacks and Defenses on Intelligent Connected Vehicles.
\item Identify the assets, actors, assumptions, and trust boundaries associated with avs, attacks, and communication.
\item Distinguish Introduction from Background and explain where their responsibilities intersect.
\item Analyze how failures in avs, attacks, and communication can affect confidentiality, integrity, availability, privacy, or safety.
\item Evaluate relevant controls through physical consequences, safety, availability, environmental dependencies, remote connectivity, maintenance, and recovery.
\item Audit an implementation using asset and dependency maps, physical-access records, sensor telemetry, safety constraints, maintenance logs, and recovery exercises.
\item Prioritize remediation according to exploitability, consequence, control strength, and recovery readiness.
\item Verify that corrective actions change observable risk rather than documentation alone.
\end{itemize}
\section*{Cornell Cue-and-Notes}
\Needspace{8\baselineskip}
\subsection*{Introduction}
\begin{longtable}{>{\RaggedRight\arraybackslash}p{0.29\textwidth}|>{\RaggedRight\arraybackslash}p{0.65\textwidth}}
\CornellHeader
\CornellNoteRow{What security problem does Introduction address?}{\textbf{Core idea.} Treat Introduction as a structured part of An Overview of Cyber Attacks and Defenses on Intelligent Connected Vehicles, with particular attention to icvs, attacks, avs, and communication. \textbf{Analytical lens.} This topic establishes a component of the chapter's argument; connect its definitions and mechanisms to a testable security objective. For icvs, attacks, and avs, model safety and physical consequences alongside conventional information-security effects. \textbf{Security reasoning.} Identify what must remain protected, who or what can alter the expected state, which assumptions cross a trust boundary, and how failure changes confidentiality, integrity, availability, privacy, or safety. \textbf{Application.} The practitioner should evaluate cyber and physical failure together, enforce safe states, and test operational recovery. \textbf{Evidence.} Seek asset and dependency maps, physical-access records, sensor telemetry, safety constraints, maintenance logs, and recovery exercises. Related subtopics include Organization.}
\end{longtable}
\begin{CornellSummaryBox}{Topic Summary}
Introduction is best remembered as the connection between icvs, attacks, avs, and communication and the chapter's larger control objective. A defensible review states the assumption, expected behavior, evidence, failure consequence, and required response.
\end{CornellSummaryBox}
\Needspace{8\baselineskip}
\subsection*{Background}
\begin{longtable}{>{\RaggedRight\arraybackslash}p{0.29\textwidth}|>{\RaggedRight\arraybackslash}p{0.65\textwidth}}
\CornellHeader
\CornellNoteRow{Which assumptions matter in Background?}{\textbf{Core idea.} Treat Background as a structured part of An Overview of Cyber Attacks and Defenses on Intelligent Connected Vehicles, with particular attention to attack, avs, potential, and assets. \textbf{Analytical lens.} This topic establishes a component of the chapter's argument; connect its definitions and mechanisms to a testable security objective. For attack, avs, and potential, model safety and physical consequences alongside conventional information-security effects. \textbf{Security reasoning.} Identify what must remain protected, who or what can alter the expected state, which assumptions cross a trust boundary, and how failure changes confidentiality, integrity, availability, privacy, or safety. \textbf{Application.} The practitioner should evaluate cyber and physical failure together, enforce safe states, and test operational recovery. \textbf{Evidence.} Seek asset and dependency maps, physical-access records, sensor telemetry, safety constraints, maintenance logs, and recovery exercises. Related subtopics include Overview, Literature Review, Attacker Model Characterization, and Threats and Vulnerabilities Enumeration.}
\end{longtable}
\begin{CornellSummaryBox}{Topic Summary}
Background is best remembered as the connection between attack, avs, potential, and assets and the chapter's larger control objective. A defensible review states the assumption, expected behavior, evidence, failure consequence, and required response.
\end{CornellSummaryBox}
\Needspace{8\baselineskip}
\subsection*{Challenges}
\begin{longtable}{>{\RaggedRight\arraybackslash}p{0.29\textwidth}|>{\RaggedRight\arraybackslash}p{0.65\textwidth}}
\CornellHeader
\CornellNoteRow{How should a reviewer evaluate Challenges?}{\textbf{Core idea.} Treat Challenges as a structured part of An Overview of Cyber Attacks and Defenses on Intelligent Connected Vehicles, with particular attention to vehicle, avs, communication, and adas. \textbf{Analytical lens.} This topic establishes a component of the chapter's argument; connect its definitions and mechanisms to a testable security objective. For vehicle, avs, and communication, model safety and physical consequences alongside conventional information-security effects. \textbf{Security reasoning.} Identify what must remain protected, who or what can alter the expected state, which assumptions cross a trust boundary, and how failure changes confidentiality, integrity, availability, privacy, or safety. \textbf{Application.} The practitioner should evaluate cyber and physical failure together, enforce safe states, and test operational recovery. \textbf{Evidence.} Seek asset and dependency maps, physical-access records, sensor telemetry, safety constraints, maintenance logs, and recovery exercises. Related subtopics include Challenges in AV Forensics, Data Verification, Autonomous Driving Assistance System, and (Adas).}
\end{longtable}
\begin{CornellSummaryBox}{Topic Summary}
Challenges is best remembered as the connection between vehicle, avs, communication, and adas and the chapter's larger control objective. A defensible review states the assumption, expected behavior, evidence, failure consequence, and required response.
\end{CornellSummaryBox}
\Needspace{8\baselineskip}
\subsection*{Av Forensics}
\begin{longtable}{>{\RaggedRight\arraybackslash}p{0.29\textwidth}|>{\RaggedRight\arraybackslash}p{0.65\textwidth}}
\CornellHeader
\CornellNoteRow{Why can Av Forensics fail in practice?}{\textbf{Core idea.} Treat Av Forensics as a structured part of An Overview of Cyber Attacks and Defenses on Intelligent Connected Vehicles, with particular attention to digital, forensic, avs, and service. \textbf{Analytical lens.} This topic is evidence-centered: define observable signals, collection points, analytic limits, escalation criteria, and preservation needs. For digital, forensic, and avs, model safety and physical consequences alongside conventional information-security effects. \textbf{Security reasoning.} Identify what must remain protected, who or what can alter the expected state, which assumptions cross a trust boundary, and how failure changes confidentiality, integrity, availability, privacy, or safety. \textbf{Application.} The practitioner should evaluate cyber and physical failure together, enforce safe states, and test operational recovery. \textbf{Evidence.} Seek asset and dependency maps, physical-access records, sensor telemetry, safety constraints, maintenance logs, and recovery exercises. Related subtopics include Safety and Reliability, Input Capture, Defense Strategies, and Attacks on AV Forensics.}
\end{longtable}
\begin{CornellSummaryBox}{Topic Summary}
Av Forensics is best remembered as the connection between digital, forensic, avs, and service and the chapter's larger control objective. A defensible review states the assumption, expected behavior, evidence, failure consequence, and required response.
\end{CornellSummaryBox}
\Needspace{8\baselineskip}
\subsection*{Communications}
\begin{longtable}{>{\RaggedRight\arraybackslash}p{0.29\textwidth}|>{\RaggedRight\arraybackslash}p{0.65\textwidth}}
\CornellHeader
\CornellNoteRow{What security problem does Communications address?}{\textbf{Core idea.} Treat Communications as a structured part of An Overview of Cyber Attacks and Defenses on Intelligent Connected Vehicles, with particular attention to malicious, bus, communication, and actors. \textbf{Analytical lens.} This topic establishes a component of the chapter's argument; connect its definitions and mechanisms to a testable security objective. For malicious, bus, and communication, model safety and physical consequences alongside conventional information-security effects. \textbf{Security reasoning.} Identify what must remain protected, who or what can alter the expected state, which assumptions cross a trust boundary, and how failure changes confidentiality, integrity, availability, privacy, or safety. \textbf{Application.} The practitioner should evaluate cyber and physical failure together, enforce safe states, and test operational recovery. \textbf{Evidence.} Seek asset and dependency maps, physical-access records, sensor telemetry, safety constraints, maintenance logs, and recovery exercises. Related subtopics include Data Manipulation, Subvert Trust Controls, Execution Guardrails, and Sensor Tampering.}
\end{longtable}
\begin{CornellSummaryBox}{Topic Summary}
Communications is best remembered as the connection between malicious, bus, communication, and actors and the chapter's larger control objective. A defensible review states the assumption, expected behavior, evidence, failure consequence, and required response.
\end{CornellSummaryBox}
\Needspace{8\baselineskip}
\subsection*{Over-The-Air Updates}
\begin{longtable}{>{\RaggedRight\arraybackslash}p{0.29\textwidth}|>{\RaggedRight\arraybackslash}p{0.65\textwidth}}
\CornellHeader
\CornellNoteRow{Which assumptions matter in Over-The-Air Updates?}{\textbf{Core idea.} Treat Over-The-Air Updates as a structured part of An Overview of Cyber Attacks and Defenses on Intelligent Connected Vehicles, with particular attention to ota, updates, software, and attacks. \textbf{Analytical lens.} This topic establishes a component of the chapter's argument; connect its definitions and mechanisms to a testable security objective. For ota, updates, and software, model safety and physical consequences alongside conventional information-security effects. \textbf{Security reasoning.} Identify what must remain protected, who or what can alter the expected state, which assumptions cross a trust boundary, and how failure changes confidentiality, integrity, availability, privacy, or safety. \textbf{Application.} The practitioner should evaluate cyber and physical failure together, enforce safe states, and test operational recovery. \textbf{Evidence.} Seek asset and dependency maps, physical-access records, sensor telemetry, safety constraints, maintenance logs, and recovery exercises. Related subtopics include Active Scanning, Attacks on OTA Updates, Internal Spearphishing, and Defense Strategies.}
\end{longtable}
\begin{CornellSummaryBox}{Topic Summary}
Over-The-Air Updates is best remembered as the connection between ota, updates, software, and attacks and the chapter's larger control objective. A defensible review states the assumption, expected behavior, evidence, failure consequence, and required response.
\end{CornellSummaryBox}
\Needspace{8\baselineskip}
\subsection*{Open Problems}
\begin{longtable}{>{\RaggedRight\arraybackslash}p{0.29\textwidth}|>{\RaggedRight\arraybackslash}p{0.65\textwidth}}
\CornellHeader
\CornellNoteRow{How should a reviewer evaluate Open Problems?}{\textbf{Core idea.} Treat Open Problems as a structured part of An Overview of Cyber Attacks and Defenses on Intelligent Connected Vehicles, with particular attention to false, attacker, true, and communication. \textbf{Analytical lens.} This topic establishes a component of the chapter's argument; connect its definitions and mechanisms to a testable security objective. For false, attacker, and true, model safety and physical consequences alongside conventional information-security effects. \textbf{Security reasoning.} Identify what must remain protected, who or what can alter the expected state, which assumptions cross a trust boundary, and how failure changes confidentiality, integrity, availability, privacy, or safety. \textbf{Application.} The practitioner should evaluate cyber and physical failure together, enforce safe states, and test operational recovery. \textbf{Evidence.} Seek asset and dependency maps, physical-access records, sensor telemetry, safety constraints, maintenance logs, and recovery exercises. Related subtopics include OTA Updates Masquerading, and Multiple Choice.}
\end{longtable}
\begin{CornellSummaryBox}{Topic Summary}
Open Problems is best remembered as the connection between false, attacker, true, and communication and the chapter's larger control objective. A defensible review states the assumption, expected behavior, evidence, failure consequence, and required response.
\end{CornellSummaryBox}
\begin{CornellWarningBox}{Historical Context}
Some products, protocols, examples, threat observations, or legal references in this chapter are historically bounded. Retain the underlying threat model and control objective, but verify present applicability before treating a named implementation as current guidance.
\end{CornellWarningBox}
\begin{CornellOverviewBox}{Defensive Guidance}
Apply the chapter by making assets, authority, trust, expected behavior, and failure response explicit. In practice, evaluate cyber and physical failure together, enforce safe states, and test operational recovery. A control is not fully supported until its operation, monitoring, ownership, and recovery behavior are demonstrated with evidence.
\end{CornellOverviewBox}
\section*{Threat-and-Control Map}
\begingroup\scriptsize\setlength{\tabcolsep}{2pt}
\begin{longtable}{>{\RaggedRight\arraybackslash}p{0.135\textwidth}>{\RaggedRight\arraybackslash}p{0.145\textwidth}>{\RaggedRight\arraybackslash}p{0.155\textwidth}>{\RaggedRight\arraybackslash}p{0.145\textwidth}>{\RaggedRight\arraybackslash}p{0.16\textwidth}>{\RaggedRight\arraybackslash}p{0.17\textwidth}}
\toprule\rowcolor{Navy}\color{white}\textbf{Asset} & \color{white}\textbf{Threat / Failure} & \color{white}\textbf{Vulnerability} & \color{white}\textbf{Impact} & \color{white}\textbf{Detection / Evidence} & \color{white}\textbf{Control}\\\midrule
\endfirsthead
\toprule\rowcolor{Navy}\color{white}\textbf{Asset} & \color{white}\textbf{Threat / Failure} & \color{white}\textbf{Vulnerability} & \color{white}\textbf{Impact} & \color{white}\textbf{Detection / Evidence} & \color{white}\textbf{Control}\\\midrule
\endhead
People, physical processes, connected devices, and essential services & Tampering, unsafe command, service disruption, surveillance, or cascading failure & Insecure remote access, fragile dependencies, weak update paths, or insufficient safety separation & Safety events, prolonged outage, physical damage, privacy harm, or public disruption & Operational telemetry, integrity monitoring, physical controls, and anomaly investigation & Layered cyber-physical protection, safe-state design, segmentation, secure maintenance, and rehearsed recovery\\\midrule
Assumptions surrounding avs & Misuse or unexpected state transition & Trust in attacks is not validated & A local weakness becomes a broader compromise path & Review evidence for communication and test negative paths & Make trust explicit, constrain authority, and fail safely\\\midrule
Recovery capability and decision evidence & Control failure remains unnoticed or cannot be reversed & Monitoring, ownership, or restoration has not been exercised & Longer exposure and slower, less reliable recovery & Alert-to-response exercises and restoration tests & Assigned ownership, actionable telemetry, preserved evidence, and rehearsed recovery\\\midrule
\bottomrule
\end{longtable}
\endgroup
\section*{Practical Security Checklist}
\begin{itemize}[label=$\square$]
\item Define the in-scope assets, users, services, data, and operational dependencies for An Overview of Cyber Attacks and Defenses on Intelligent Connected Vehicles.
\item Document the expected behavior and security objective of avs.
\item Map identities, privileges, trust boundaries, and externally controlled inputs associated with avs and attacks.
\item Record the threat or failure being addressed, its prerequisites, likely path, and credible consequences.
\item Inspect asset and dependency maps, physical-access records, sensor telemetry, safety constraints, maintenance logs, and recovery exercises.
\item Test both the intended path and negative paths, including invalid state, partial failure, excessive privilege, and unavailable dependencies.
\item Confirm preventive controls are complemented by actionable detection and a named response owner.
\item Exercise containment, restoration, and communication; record actual recovery behavior and unresolved dependencies.
\item Validate exceptions, compensating controls, expiration dates, and accountable approval.
\item Preserve reproducible evidence and tie each finding to affected assets, risk, remediation, and verification criteria.
\end{itemize}
\begin{CornellSummaryBox}{Chapter Synthesis}
The chapter's principal lesson is that an overview of cyber attacks and defenses on intelligent connected vehicles must be evaluated as an operating security system rather than an isolated product or statement. The most important failure patterns involve avs, attacks, communication, and ota, especially when ownership, boundaries, telemetry, or recovery remain implicit. A reviewer should connect architecture and behavior to asset and dependency maps, physical-access records, sensor telemetry, safety constraints, maintenance logs, and recovery exercises, prioritize findings by reachable consequence and control independence, and verify remediation through observable change. Within Connected and Automated Vehicle Security, this chapter contributes a reusable method: state the objective, expose assumptions, test failure, preserve evidence, and learn from operation.
\end{CornellSummaryBox}
\section*{Self-Test Questions}
\begin{enumerate}
\item What is the central security objective of An Overview of Cyber Attacks and Defenses on Intelligent Connected Vehicles?
\item How does Introduction contribute to the chapter's broader security argument?
\item Distinguish Background from Challenges.
\item Which trust assumptions should be made explicit when evaluating avs?
\item How could a weakness involving attacks affect confidentiality, integrity, availability, privacy, or safety?
\item What evidence would demonstrate that controls surrounding communication operate as intended?
\item Why should prevention, detection, response, and recovery be evaluated together in this domain?
\item Scenario: A connected physical process remains functionally available after a cyber event but no longer operates within safe constraints. What should the reviewer examine first, and why?
\item Scenario: A control associated with ota passes a configuration check but fails during an operational exercise. How should the finding be interpreted and prioritized?
\item How should a practitioner distinguish enduring principles in this chapter from time-sensitive tools, examples, or implementation details?
\end{enumerate}
\Needspace{8\baselineskip}
\section*{Answer Key}
\begin{enumerate}
\item The objective is to protect the relevant assets by understanding physical consequences, safety, availability, environmental dependencies, remote connectivity, maintenance, and recovery and by connecting controls to measurable risk reduction.
\item Introduction provides one part of the causal chain: it should identify expected behavior, dependencies, failure modes, and the evidence needed to justify confidence.
\item Background and Challenges address different portions of the problem. A strong comparison states their scope, assumptions, enforcement points, evidence, and how a failure in one affects the other.
\item Reviewers should identify who controls avs, what inputs or dependencies it trusts, where privilege changes, what happens during partial failure, and how those assumptions are tested.
\item A weakness in attacks can propagate across trust boundaries. The actual impact depends on reachable assets, granted authority, control independence, visibility, and recovery capability.
\item Useful evidence includes asset and dependency maps, physical-access records, sensor telemetry, safety constraints, maintenance logs, and recovery exercises; configuration alone is insufficient when operation, monitoring, or recovery is part of the claim.
\item Prevention reduces probability, detection limits dwell time, response limits spread, and recovery restores objectives. Omitting one layer creates an unexamined failure mode.
\item Begin by establishing scope, ownership, expected behavior, and the violated assumption. Then evaluate cyber and physical failure together, enforce safe states, and test operational recovery, preserving evidence for prioritization and follow-up.
\item Treat the exercise as stronger evidence than a static check. Prioritize according to reachable impact, likelihood of real operating conditions, missing compensating controls, and recovery difficulty.
\item Retain the principle, threat model, and control objective; separately label technologies, legal references, product examples, and threat observations whose applicability depends on time or environment.
\end{enumerate}
\section*{Key-Term Recap}
\begin{longtable}{>{\RaggedRight\arraybackslash}p{0.27\textwidth}>{\RaggedRight\arraybackslash}p{0.66\textwidth}}
\toprule\rowcolor{Navy}\color{white}\textbf{Term} & \color{white}\textbf{Meaning}\\\midrule
\endfirsthead
\toprule\rowcolor{Navy}\color{white}\textbf{Term} & \color{white}\textbf{Meaning}\\\midrule
\endhead
\textbf{Threat} & A circumstance or actor capable of causing harm by exploiting a weakness or adverse condition. \\\midrule
\textbf{Forensics} & The use of repeatable methods to preserve and analyze evidence for investigative or legal purposes. \\\midrule
\textbf{Risk} & The effect of uncertainty on objectives, commonly evaluated through likelihood, impact, exposure, and context. \\\midrule
\textbf{Asset} & Information, service, capability, or physical resource whose value justifies protection. \\\midrule
\textbf{Threat Model} & A structured representation of assets, trust boundaries, adversaries, attack paths, and mitigations. \\\midrule
\textbf{Integrity} & The property that information and systems remain accurate, complete, and protected from unauthorized modification. \\\midrule
\textbf{Authentication} & The process of establishing confidence that an asserted identity is genuine. \\\midrule
\textbf{Intrusion Detection} & Monitoring and analysis intended to identify unauthorized or malicious activity. \\\midrule
\textbf{Chain Of Custody} & The documented history of evidence possession, handling, transfer, and integrity. \\\midrule
\textbf{Encryption} & A keyed transformation of plaintext into ciphertext to protect confidentiality. \\\midrule
\bottomrule
\end{longtable}
\begin{CornellExamBox}{One-Sentence Takeaway}
Treat an overview of cyber attacks and defenses on intelligent connected vehicles as a verifiable chain from assets and assumptions through controls, evidence, response, and recovery.
\end{CornellExamBox}
\end{document}
